% =====================================================================
%  tikzlib · ciencias/orbitales — diagramas de orbitales (cajas + flechas)
%  Requiere: \usepackage{tikz}  (+ amsmath para las flechas)
%  Flechas:  \flup (↑)  \fldn (↓)  \fludn (↑↓)
%  \sub{etiqueta}{n.º de orbitales}        → cajas VACÍAS (p. ej. \sub{2p}{3})
%  \subll{etiqueta}{lista}                 → cajas LLENAS
%        p. ej.  \subll{2p}{\fludn,\flup,\flup}
%  Preview: tikzlib/previews/orbitales.png
% =====================================================================
\newcommand{\flup}{\ensuremath{\uparrow}}
\newcommand{\fldn}{\ensuremath{\downarrow}}
\newcommand{\fludn}{\ensuremath{\uparrow\!\downarrow}}

\newcommand{\sub}[2]{%
\begin{tikzpicture}[baseline=-0.6ex]
  \foreach \i in {1,...,#2}{%
    \draw[line width=0.6pt] ({(\i-1)*0.72},0) rectangle ({(\i-1)*0.72+0.66},0.66);}
  \node[font=\sffamily\footnotesize\bfseries] at ({#2*0.72/2-0.03},-0.28) {#1};
\end{tikzpicture}}

\newcommand{\subll}[2]{%
\begin{tikzpicture}[baseline=-0.6ex]
  \foreach \c [count=\j from 0] in {#2}{%
    \draw[line width=0.6pt] ({\j*0.72},0) rectangle ({\j*0.72+0.66},0.66);
    \node at ({\j*0.72+0.33},0.33) {\c};
    \xdef\last{\j}}
  \node[font=\sffamily\footnotesize\bfseries] at ({(\last+1)*0.72/2-0.03},-0.28) {#1};
\end{tikzpicture}}
